\documentclass[11pt,a4paper]{ltjsarticle}

\usepackage{amsmath,amssymb,amsthm}
\usepackage{geometry}
\usepackage{algorithm}
\usepackage{algpseudocode}
\usepackage{enumitem}
\usepackage{hyperref}

\geometry{a4paper,margin=25mm}

\theoremstyle{plain}
\newtheorem{theorem}{定理}[section]
\newtheorem{proposition}[theorem]{命題}
\newtheorem{lemma}[theorem]{補題}
\newtheorem{corollary}[theorem]{系}
\theoremstyle{definition}
\newtheorem{definition}[theorem]{定義}
\theoremstyle{remark}
\newtheorem{remark}[theorem]{注意}
\renewcommand{\proofname}{証明}

\newcommand{\R}{\mathbb{R}}

\title{上下限制約付き双対単体法における自明な初期基底の構成\\
  ---非有界変数の $M$ による有界化とゼロ固定スラック変数による統一的手法---}
\author{}
\date{\today}

\begin{document}
\maketitle

\begin{abstract}
単体法は，線型計画問題を解くためのよく知られた枠組みであるが，一般の線形計画問題に対しては，ネットワークフロー問題や右辺が非負な不等式標準形など一部の特殊なクラスを除き，実行可能な初期基底解を計算なしに構成する一般的な方法は知られておらず，実務上は2段階法やBig$-M$法，あるいはクラッシュ手続きのような補助的な手続きを要する．一方，主実行可能性ではなく\emph{双対実行可能性}に着目すると，コスト係数の符号のみから非基底変数の配置(下限か上限か)を決定でき，全スラック基底は計算なしに双対実行可能な初期基底となる．ただし，この構成は非基底変数を有限な上下限に配置できることを前提とするため，下限・上限が非有界な変数に対しては直接適用できない．任意の線形計画問題に対して，非有界な下限・上限を任意の正数 $M$ を用いてそれぞれ $-M,+M$ に置き換えて問題をあらかじめ有界化した上では，各等式制約に値を $0$ に固定したスラック変数を追加することにより，フェーズ1を要さずに自明な双対実行可能初期基底を構成し，自明な基底を用いて上下限制約付き双対単体法をただちに開始できる．本稿では，その具体的なアルゴリズムの実際の計算式を，任意に大きい$M$に対して解析することで，非有界な変数を持つ一般の線型計画問題に対して初期基底の構築が不要な単体法の拡張手法を与える．
\end{abstract}

\section{はじめに}

線形計画問題(Linear Programming, LP)に対する単体法は，理論と実装の両面で長年研究されてきた古典的な解法である．単体法を実際に動かすには，制約行列 $A$ から可逆な $m\times m$ 部分行列(基底)$A_B$ を選び，非基底変数をその上限または下限に固定したときに基底変数の値
\[
  x_B = A_B^{-1}\bigl(b - A_N x_N\bigr)
\]
がすべて自身の上下限に収まる，すなわち\emph{実行可能な}初期基底解から出発する必要がある．

この初期基底の構築は，一般には自明ではない．例えば，制約がすべて「$\le$」形であり右辺 $b$ が非負であるような標準形では，スラック変数を追加した全スラック基底がそのまま実行可能基底となる．また，輸送問題や割当問題などのネットワークフロー問題では，供給量・需要量の整合性のもとで全域木構造から実行可能基底を組合せ的に構成できることが知られている．しかし，これらは制約行列やデータに特別な構造を仮定した特殊なクラスであり，等式制約 $Ax=b$ と一般の範囲制約 $l\le x\le u(l,u \in \mathbb R \cup \{-\infty, +\infty\})$ からなる\emph{任意の}線形計画問題に対して，計算コストをかけずに実行可能な初期基底を与える一般的な方法は知られていない．実際，初期実行可能基底の存在判定およびその構成は，それ自体が一つの線形計画問題(補助問題)を解くことと本質的に同程度の難しさを持ち，このために2段階法(補助目的関数を最小化するフェーズ1を別途実行する)や大$M$法(人工変数に大きなペナルティ係数を課して単一の目的関数で最適化する)といった補助的な手続きが用いられてきた．これらはいずれも，追加の変数や補助問題の構築・数値的なパラメータ選択(大$M$法における $M$ の大きさ)を要する点で，「計算なしに得られる自明な初期基底」からは程遠い．

本稿では，有界な上下限制約付き線型計画問題は自明な初期基底が構築可能なことに着目し，非有界変数の上下限$\pm \infty$を一旦有界な$\pm M$に置き換えた場合の双対単体法のアルゴリズムを解析し、拡張手法を構築することでこの困難を回避する．上下限制約付き単体法において，非基底変数 $x_j$ がその下限 $l_j$ にあるか上限 $u_j$ にあるかは，被約費用(reduced cost) $\bar c_j$ の符号のみによって双対実行可能性が定まる．とりわけ，コスト $0$ のスラック変数のみからなる全スラック基底では，被約費用は $\bar c_j=c_j$ そのものに一致するため，非基底の原変数を
\[
  c_j\ge 0 \ \Longrightarrow\ x_j=l_j,\qquad
  c_j< 0 \ \Longrightarrow\ x_j=u_j
\]
と配置するだけで，\emph{計算なしに}双対実行可能な初期基底が得られる．この自明な初期基底を使用可能にするように単体法を拡張するのである。

本稿の構成は次の通りである．第2節で上下限制約付き単体法および双対単体法の基礎を整理する．第3節で，主実行可能な初期基底の構築が一般に困難である事情と，それが成立する特殊なクラスを確認する．第4節で，非有界変数を $M$ で有界化した上でのゼロ固定スラック変数による自明な初期基底の構成と，それに続く上下限制約付き双対単体法のアルゴリズムを，実際の計算式とともに具体的に記述する．さらに，$M\to+\infty$ での挙動を解析することにより，初期基底の双対実行可能性，$M$ を十分大きく取ることの正当性，および元問題の実行不可能性・非有界性・有限最適解の存在を判別できることまでを一貫して示す．退化に対する巡回防止規則や，効率的な基底更新手法との統合といった実装上の課題は今後の検討課題とする．

\section{準備：上下限制約付き単体法}

\subsection{問題設定}

本稿では次の形の線形計画問題を扱う．
\begin{align}
  \min\quad & c^\top x \label{eq:lp}\\
  \text{s.t.}\quad & Ax = b, \notag\\
  & l_j \le x_j \le u_j, \quad j=1,\dots,n, \notag
\end{align}
ここで $A\in\R^{m\times n}$，$b\in\R^m$，$c\in\R^n$ であり，$\operatorname{rank}(A)=m$ とする(行の一次従属性は前処理で除去済みと仮定する)．各変数の上下限 $l_j,u_j$ は
\[
  l_j \in \R\cup\{-\infty\},\qquad u_j \in \R\cup\{+\infty\},\qquad l_j\le u_j
\]
を満たす拡大実数値をとってよい．この定式化は，$l=0,\,u=(+\infty,\dots,+\infty)$ とすれば通常の標準形を，$l=-\infty,\,u=+\infty$ とすれば自由変数を含む問題を包含する，最も一般的な設定である．

\subsection{有界な上下限制約付き問題の基底解・被約費用・実行可能性}

添字集合 $\{1,\dots,n\}$ を基底添字集合 $B$($|B|=m$，$A_B$ は正則)と非基底添字集合 $N=\{1,\dots,n\}\setminus B$ に分割する．各非基底変数には状態 $\eta(j)\in\{L,U\}$ を割り当て，
\[
  x_j =
  \begin{cases}
    l_j & (\eta(j)=L)\\
    u_j & (\eta(j)=U)
  \end{cases}
  \qquad (j\in N)
\]
とする．このとき基底変数は
\[
  x_B = A_B^{-1}\bigl(b - A_N x_N\bigr)
\]
により一意に定まる．組 $(B,\eta,x)$ が\emph{主実行可能}であるとは，$l_B\le x_B\le u_B$ が成り立つことをいう．

双対変数(単体乗数) $y^\top = c_B^\top A_B^{-1}$ を用いて，各変数 $j$ の被約費用を
\[
  \bar c_j = c_j - y^\top A_j
\]
と定める．基底変数については常に $\bar c_j=0$ である．組 $(B,\eta)$ が\emph{双対実行可能}であるとは，
\[
  j\in N,\ \eta(j)=L \ \Longrightarrow\ \bar c_j\ge 0,\qquad
  j\in N,\ \eta(j)=U \ \Longrightarrow\ \bar c_j\le 0
\]
が成り立つことをいう．主実行可能性と双対実行可能性が同時に成り立つ基底は，LPの双対定理により\emph{最適基底}である．

\subsection{上下限制約付き双対単体法}

双対単体法は，双対実行可能な基底から出発し，双対実行可能性を保ったまま主実行不可能性を解消していく反復法であり，主実行可能性に到達した時点でそれが最適解であることが保証される．1回の反復は次のように行う．

\begin{enumerate}[label=(\roman*)]
\item \textbf{脱退変数の選択}：$x_B$ のうち上下限を逸脱している行 $r$ を一つ選ぶ．逸脱量が最大の行を選ぶのが典型的である(退化に対しては Bland 則等の巡回防止規則を援用する)．逸脱の方向を
  \[
    d =
    \begin{cases}
      +1 & (x_{B,r} < l_{\beta(r)}, \text{増加させたい})\\
      -1 & (x_{B,r} > u_{\beta(r)}, \text{減少させたい})
    \end{cases}
  \]
  とおく．ここで $\beta(r)$ は行 $r$ に対応する基底変数の添字である．
\item \textbf{ピボット行の計算}：$\alpha^\top := e_r^\top A_B^{-1}A_N$(タブローの行 $r$ の非基底部分)を計算する．各非基底 $j$ に対し符号 $\sigma_j:=+1\,(\eta(j)=L)$，$\sigma_j:=-1\,(\eta(j)=U)$ を定め，$\hat\alpha_j:=\sigma_j\alpha_j$，$\hat c_j:=\sigma_j\bar c_j\ (\ge 0,\text{双対実行可能性より})$ とおく．
\item \textbf{加入可能集合と比検定}：加入可能な非基底変数の集合を
  \[
    \mathrm{Eligible} := \{\, j\in N : d\,\hat\alpha_j < 0 \,\}
  \]
  と定める．$\mathrm{Eligible}=\emptyset$ ならば，現在の問題(の上下限制約のもとで)は実行不可能であり，アルゴリズムは停止する．そうでなければ，
  \[
    q \in \operatorname*{arg\,min}_{j\in \mathrm{Eligible}} \frac{\hat c_j}{|\hat\alpha_j|}
  \]
  を加入変数として選ぶ．
\item \textbf{ピボット}：$\beta(r)$ を非基底化し，$d=+1$ ならば $l_{\beta(r)}$ に，$d=-1$ ならば $u_{\beta(r)}$ に固定する．$q$ を基底化し，基底 $B$，逆行列 $A_B^{-1}$，基底解 $x_B$，非基底の状態 $\eta$ を更新する．
\end{enumerate}

この手続きは，各反復で双対実行可能性を保ったまま，退化がなければ双対目的関数を真に改善するため，有限回の反復で「すべての行が上下限を満たす($x_B$ が実行可能)」か「$\mathrm{Eligible}=\emptyset$(実行不可能)」のいずれかに到達する．前者に至れば，得られた $(B,\eta,x)$ は主・双対ともに実行可能であるから，問題~\eqref{eq:lp}の最適解である．

\section{初期基底構築の困難性}

\subsection{自明な主実行可能基底を持つ特殊なクラス}

上下限制約付き単体法(主単体法)を実行するには，本来，主実行可能な初期基底が必要である．次のような特殊なクラスでは，これを計算なしに構成できる．

\begin{itemize}
\item \textbf{右辺非負の不等式標準形}：制約が $Ax\le b,\ b\ge 0,\ x\ge 0$ の形であれば，スラック変数 $s=b-Ax\ge 0$ を導入した全スラック基底 $A_B=I$，$x_N=0$ がそのまま主実行可能である．
\item \textbf{ネットワークフロー問題}：輸送問題・割当問題・最小費用流問題などでは，制約行列が(全域木に対応する)完全単模性を持ち，供給量・需要量の整合性のもとで，全域木構造から組合せ的に実行可能な初期基底を構成できる(ネットワーク単体法)．
\end{itemize}

これらはいずれも，制約行列や右辺ベクトルに特別な構造を仮定している点で特殊である．

\subsection{自明な双対実行可能基底を持つ特殊なクラス}
\label{sec:dual-feasible-classes}

主実行可能性ではなく双対実行可能性に着目すると，事情は異なる．第2.2節で見たように，コスト $0$ のスラック変数のみからなる全スラック基底 $B_0$(すなわち $A_{B_0}=I$)では，単体乗数が自明に $y=0$ となるため，非基底の原変数 $j$ の被約費用は $\bar c_j=c_j$ に一致する．したがって，非基底変数を
\[
  c_j\ge 0 \ \Longrightarrow\ \eta(j)=L,\qquad
  c_j< 0 \ \Longrightarrow\ \eta(j)=U
\]
と配置しさえすれば，計算(基底逆行列や比検定)を一切行うことなく双対実行可能性の条件が満たされる．ただし，この配置が\emph{定義できる}ためには，配置先となる側の上下限が有限でなければならない．この点に注意すると，次のようなクラスでは，主実行可能性の成否とは無関係に，全スラック基底が自明に双対実行可能となる．

\begin{itemize}
\item \textbf{有界変数(箱型制約)問題}：すべての変数が両側とも有限な上下限 $l_j,u_j\in\R$ を持つ場合(bounded-variable / box-constrained LP)，コスト $c_j$ の符号がどちらであっても，配置すべき側に必ず有限な上下限が存在する．したがって全スラック基底は，$c$ の値によらず常に双対実行可能である．
\item \textbf{標準形かつコスト非負}：$l=0,\,u=+\infty$(通常の標準形)であって，かつ $c\ge0$ が成り立つ場合，すべての非基底変数を下限 $0$ に配置すれば双対実行可能性が満たされる．とくに，第3.1節の「右辺非負の不等式標準形」の条件($b\ge0$)も同時に成り立つならば，全スラック基底は主・双対ともに実行可能であり，反復を一切行わずに $x=0$ が直ちに最適解であると判定できる．
\end{itemize}

これらのクラスに共通するのは，各変数のコスト $c_j$ の符号が，その変数の少なくとも一方の\emph{有限な}上下限と整合している，という点である．逆に，この構成が破綻するのは，ある変数 $j$ について，コストの符号が指し示す側の上下限が非有界である場合，すなわち $c_j\ge0$ にもかかわらず $l_j=-\infty$ であるか，$c_j<0$ にもかかわらず $u_j=+\infty$ である場合に限られる．一般の線形計画問題ではこのような変数の存在を排除できないため，双対実行可能性からの出発もまた，これらの特殊なクラスを除けば自明ではない．

\subsection{一般の等式制約付き問題における困難性}

一般の等式制約 $Ax=b$ と上下限制約 $l\le x\le u$ からなる問題に対しては，主実行可能性の側にも双対実行可能性の側にも，前2節で見たような構造上の手がかりがない．主実行可能性については，与えられた基底 $B$ に対して主実行可能な非基底状態 $\eta$ が存在するかどうかを判定し，存在するならそれを見つけるという問題が，それ自体「$\{x: Ax=b,\ l\le x\le u\}\neq\emptyset$ か」を判定する実行可能性問題(補助LP)と本質的に同等の難しさを持つ．このため実務上は次のいずれかの補助的手続きが用いられる．

\begin{itemize}
\item \textbf{2段階法}：人工変数 $a\ge 0$ を導入した補助問題 $\min \mathbf 1^\top a$ を(全スラック/人工変数基底から)フェーズ1として解き，実行可能基底を得てからフェーズ2で本来の目的関数を最適化する．
\item \textbf{大$M$法}：人工変数 $a\ge0$ に大きな係数 $M$ を課した合成目的関数 $c^\top x + M\mathbf 1^\top a$ を単一フェーズで最小化する．
\item \textbf{クラッシュ手続き}：制約行列 $A$ の構造(対角優位性や三角構造への近さなど)に基づき，ヒューリスティックに(非零成分の少ない列を優先するなどの規則で)正則な開始基底を選び出す．計算コストは小さいが，得られる基底が主実行可能である保証はなく，一般には後続のフェーズ1の反復回数を削減する前処理として用いられるに過ぎない．
\end{itemize}

いずれの方法も，補助変数の追加や十分大きい $M$ の選択という数値的に微妙な判断，あるいはヒューリスティックな規則に基づく非自明な計算を要するという点で，「計算なしに得られる自明な初期基底」からは程遠い．とりわけ大$M$法は，$M$ の取り方次第で数値的な不安定性を招くことが古くから指摘されており，実装上は2段階法やクラッシュ手続きとの併用が好まれる傾向にある．

一方，双対実行可能性については，第\ref{sec:dual-feasible-classes}節で見た構成の破綻要因(コストの符号が指し示す側の上下限が非有界であること)は，主実行可能性の場合とは異なり，たった1点にまで単純化されている．本稿は，この唯一の障害を $M$ による有界化によって除去したアルゴリズムを任意に大きい$M$に対して調べることで、解析，全スラック基底を任意の線形計画問題に対して自明に双対実行可能とする方法を，次節で具体的に示す．

\section{提案手法の導出}

\subsection{前処理：片側のみ非有界な変数の平行移動}
\label{sec:shift-preprocess}

以降の議論を簡潔にするため，あらかじめ次の前処理を行っておく．各 $j=1,\dots,n$ に対し，シフト量を
\[
  s_j \ :=\
  \begin{cases}
    u_j & (l_j=-\infty,\ u_j<+\infty)\\
    l_j & (l_j>-\infty,\ u_j=+\infty)\\
    0 & (\text{それ以外，すなわち両側とも有限，または両側とも非有界})
  \end{cases}
\]
と定め，変数を $x_j = x_j' + s_j$ と置き換える．このとき制約は $Ax' = b-As =: b'$ となり，上下限は $l_j':=l_j-s_j$，$u_j':=u_j-s_j$ となる．とくに，片側のみが非有界であった $j$ については，有限であった側の上下限がちょうど $0$ になる(たとえば $l_j=-\infty,u_j<\infty$ なら $l_j'=-\infty,u_j'=0$；$l_j>-\infty,u_j=+\infty$ なら $l_j'=0,u_j'=+\infty$)．目的関数は $c^\top x = c^\top x' + c^\top s$ となり，$c^\top s$ は $x'$ に関して定数であるから，$x'$ についての最小化は元の問題と同一の最適基底構造を与え，最適値は定数 $c^\top s$ だけ異なる(元問題の最適値は，シフト後の問題の最適値に $c^\top s$ を加えたものである)．

以下では，この前処理が済んでいるものとして，記号 $A,b,c,l,u$ を $A,b',c,l',u'$ の意味で用いる．すなわち，片側のみ非有界な変数については，有限な側の上下限が $0$ であると仮定してよい．この仮定のもとでは，非基底変数 $j$ が(元々非有界であった側とは反対の)有限な上下限に配置されている場合，その値は $M$ に依存せず常に $\hat l_j=0$ または $\hat u_j=0$ である．

\subsection{非有界変数の $M$ による有界化}

任意に $M>0$ を固定する．各 $j=1,\dots,n$ に対し，有界化された上下限を
\begin{equation}
  \hat l_j :=
  \begin{cases}
    l_j & (l_j > -\infty)\\
    -M & (l_j = -\infty)
  \end{cases}
  ,\qquad
  \hat u_j :=
  \begin{cases}
    u_j & (u_j < +\infty)\\
    M & (u_j = +\infty)
  \end{cases}
  \label{eq:trunc}
\end{equation}
と定める．元々有限であった上下限はそのまま用い，非有界であった側のみ $\mp M$ に置き換える．$\hat l_j,\hat u_j$ はすべて有限であるから，問題
\begin{equation}
  \min\ c^\top x \quad \text{s.t.}\quad Ax=b,\ \hat l\le x\le \hat u
  \label{eq:lpM}
\end{equation}
は上下限制約付き単体法がそのまま適用できる有界な問題となる．$\hat l\le l$ 側は不変，$\hat u\ge u$ 側も不変であり，$[\hat l,\hat u]\subseteq[l,u]$ であるから，\eqref{eq:lpM}の実行可能領域は元の問題\eqref{eq:lp}の実行可能領域の部分集合である．

\subsection{ゼロ固定スラック変数による自明な初期基底}

等式制約 $Ax=b$ に，各行に対応するスラック変数 $s\in\R^m$ を追加し，
\[
  Ax + s = b,\qquad l_s = u_s = 0 \ \ (\text{すなわち } s\equiv 0 \text{ に固定})
\]
と書き換える．拡大された変数ベクトルは $(x,s)\in\R^{n+m}$ であり，拡大制約行列は $\hat A = [\,A \mid I_m\,]$ である．$s$ のコストは $0$ とする．

このとき，$B_0:=\{n+1,\dots,n+m\}$(スラック変数の添字集合)を初期基底として採用すると，$A_{B_0}=I_m$ であるから，
\[
  A_{B_0}^{-1} = I_m
\]
であり，逆行列の計算すら不要である．非基底(原変数)の状態は，コスト係数の符号のみから
\begin{equation}
  \eta(j) =
  \begin{cases}
    L, \ x_j=\hat l_j & (c_j \ge 0)\\
    U, \ x_j=\hat u_j & (c_j < 0)
  \end{cases}
  \qquad (j=1,\dots,n)
  \label{eq:init-status}
\end{equation}
と定める．このとき基底(スラック)変数の値は
\begin{equation}
  x_{B_0} = s = b - A x_N
  \label{eq:init-basic}
\end{equation}
により直接計算される．$s$ は一般には $0$ に一致せず，\eqref{eq:init-basic}の値は $s$ 自身の上下限 $[0,0]$ を逸脱しうるが，これは以下に見るように双対実行可能性の観点からは問題にならない．

\begin{proposition}[初期基底の双対実行可能性]
\label{prop:dual-feasible-init}
\eqref{eq:init-status}で定めた $(B_0,\eta)$ は，問題~\eqref{eq:lpM}に対して双対実行可能である．
\end{proposition}

\begin{proof}
基底コストは $c_{B_0}=0$($s$ のコストが $0$)であり，$A_{B_0}=I$ であるから，単体乗数は $y^\top = c_{B_0}^\top A_{B_0}^{-1} = 0$ である．したがって非基底の原変数 $j$ の被約費用は
\[
  \bar c_j = c_j - y^\top A_j = c_j
\]
となり，コストそのものに一致する．\eqref{eq:init-status}より，$c_j\ge0$ のとき $\eta(j)=L$ かつ $\bar c_j=c_j\ge0$，$c_j<0$ のとき $\eta(j)=U$ かつ $\bar c_j=c_j<0\le0$ であるから，双対実行可能性の条件を満たす．基底変数($s$)の被約費用は定義より恒等的に $0$ であり，条件を課す必要がない．
\end{proof}

命題~\ref{prop:dual-feasible-init}により，$(B_0,\eta,x)$ は\emph{計算コストなしに}双対実行可能な初期基底となる．ここで本質的に用いているのは，(a) $s$ のコストが $0$ であることと $A_{B_0}=I$ であることから単体乗数が自明に $y=0$ となること，(b) 非基底変数を配置すべき上下限 $\hat l_j,\hat u_j$ がいずれも有限であること，の2点である．(b)はまさに第4.2節の $M$ による有界化によって保証されている．元の問題に非有界な変数が存在する限り，この有界化なしには\eqref{eq:init-status}の配置(特に $u_j=+\infty$ の変数を $\eta(j)=U$ に置くこと，あるいは $l_j=-\infty$ の変数を $\eta(j)=L$ に置くこと)が定義できない点に注意されたい．

\subsection{アルゴリズム}

以上から，変数を $M$ を用いて有界にした場合のアルゴリズムは以下のようになる．

\textbf{Step~1(初期基底の構成)．} 式~\eqref{eq:init-status}，\eqref{eq:init-basic}により，基底を $B\gets B_0=\{n+1,\dots,n+m\}$(スラック変数の添字)，非基底の原変数の状態を $\eta$，非基底値を $x_N$，基底値を $x_B=s=b-Ax_N$ と定める．命題~\ref{prop:dual-feasible-init}により，この $(B,\eta)$ は計算なしに双対実行可能である．

\textbf{Step~2(反復)．} 以下の(a)--(d)を，(a)で終了と判定されるまで繰り返す．

\begin{enumerate}[label=(\alph*)]
\item \textbf{主実行可能性の判定(終了条件)．} 現在の基底 $B$ に対し，基底変数の値は
\[
  x_B \ =\ A_B^{-1}\bigl(b - A_N x_N\bigr)
\]
で与えられる．ここで非基底値 $x_N$ の各成分 $x_j\,(j\in N)$ は，非基底の状態 $\eta(j)\in\{L,U\}$ に応じて
\[
  x_j =
  \begin{cases}
    \hat l_j & (\eta(j)=L)\\
    \hat u_j & (\eta(j)=U)
  \end{cases}
\]
であり，$\hat l_j,\hat u_j$ 自体も式~\eqref{eq:trunc}により，$j$ が元々非有界であったか否かで場合分けして
\[
  \hat l_j =
  \begin{cases}
    l_j & (l_j > -\infty)\\
    -M & (l_j = -\infty)
  \end{cases}
  ,\qquad
  \hat u_j =
  \begin{cases}
    u_j & (u_j < +\infty)\\
    M & (u_j = +\infty)
  \end{cases}
\]
と与えられる．すなわち，$j$ が元々下限 $-\infty$ を持ち $\eta(j)=L$ であれば $x_j=-M$，元々上限 $+\infty$ を持ち $\eta(j)=U$ であれば $x_j=M$ が，$x_B$ の計算に直接用いられることになる(初回はStep~1の式~\eqref{eq:init-basic}，すなわち $A_B=I$ の場合に一致する．2回目以降は，この式を毎回計算し直す必要はなく，後述(d)のピボット操作によって $x_B$ を逐次更新すればよい)．ここで $\beta(i)$ を行 $i$ に対応する基底変数の添字とすると，判定に用いる上下限 $\hat l_{\beta(i)},\hat u_{\beta(i)}$ は，式~\eqref{eq:trunc}により，$\beta(i)$ が元々非有界であったか否かで場合分けして
\[
  \hat l_{\beta(i)} =
  \begin{cases}
    l_{\beta(i)} & (l_{\beta(i)} > -\infty)\\
    -M & (l_{\beta(i)} = -\infty)
  \end{cases}
  ,\qquad
  \hat u_{\beta(i)} =
  \begin{cases}
    u_{\beta(i)} & (u_{\beta(i)} < +\infty)\\
    M & (u_{\beta(i)} = +\infty)
  \end{cases}
\]
と与えられる．これを用いて，$x_B$ のうち自身の上下限を逸脱している行の添字集合(逸脱制約の添字集合)を
\[
  V \ :=\ \bigl\{\, i\in\{1,\dots,m\} \ :\ x_{B,i} < \hat l_{\beta(i)} \ \text{または}\ x_{B,i} > \hat u_{\beta(i)} \,\bigr\}
\]
と定める．とくに $\beta(i)$ が元々下限 $-\infty$ を持っていた場合，$x_{B,i}<\hat l_{\beta(i)}$ は $x_{B,i}<-M$ を意味し，元々上限 $+\infty$ を持っていた場合，$x_{B,i}>\hat u_{\beta(i)}$ は $x_{B,i}>M$ を意味する．$V=\emptyset$ ならば，$x_B$ はすべての上下限を満たしており($\hat l_B\le x_B\le \hat u_B$)，すなわち $(B,\eta,x)$ は問題~\eqref{eq:lpM}に対して主実行可能である．一方，双対実行可能性はStep~1で成立し，かつ以下の(b)--(d)の各反復によって保たれる(第2.3節)ため，この時点で $(B,\eta,x)$ は主・双対ともに実行可能であり，LPの弱双対定理より問題~\eqref{eq:lpM}の最適解である．アルゴリズムを終了する．$V\neq\emptyset$ の場合は(b)に進む．

\item \textbf{脱退変数の選択．} まず，各 $i\in V$ に対し逸脱量を
\[
  \Delta_i \ :=\ \max\bigl(\hat l_{\beta(i)}-x_{B,i},\ x_{B,i}-\hat u_{\beta(i)}\bigr) \ (>0)
\]
とおく．脱退行 $r\in V$ の選び方には複数の規則があり，いずれも第4.5節で改めて扱う．

\begin{itemize}
\item \textbf{Dantzig則(最大逸脱則)}：
\[
  r \ \in\ \operatorname*{arg\,max}_{i\in V} \Delta_i .
\]
\item \textbf{最急辺規則(dual steepest edge)}：現在の基底に対するタブローの行ノルム
\[
  \gamma_i \ :=\ \sum_{j\in N} \alpha_{ij}^2, \qquad \alpha_{ij} := \bigl(A_B^{-1}A_N\bigr)_{ij}
\]
を重みとして，逸脱量を正規化した
\[
  r \ \in\ \operatorname*{arg\,max}_{i\in V} \frac{\Delta_i^2}{\gamma_i}
\]
により選ぶ．$\gamma_i$ は基底の幾何(条件の悪さ)を反映するため，Dantzig則より反復回数が少なくなる傾向があることが経験的に知られている．
\item \textbf{Devex規則}：$\gamma_i$ の代わりに，各行について $1$ で初期化した近似重み $w_i$ を，ピボットのたびに(そのピボット行の要素 $\alpha_{rj}$ のみを用いて)安価に更新しながら
\[
  r \ \in\ \operatorname*{arg\,max}_{i\in V} \frac{\Delta_i^2}{w_i}
\]
により選ぶ．$\gamma_i$ の再計算(全行にわたる $O(m)$ 個の内積計算)を避けつつ最急辺規則を近似する，実装上広く用いられる方式である．
\end{itemize}

いずれの規則で $r$ を選んだ場合も，逸脱の方向は
\[
  d \ :=\
  \begin{cases}
    +1 & (x_{B,r} < \hat l_{\beta(r)}, \ \text{増加させたい})\\
    -1 & (x_{B,r} > \hat u_{\beta(r)}, \ \text{減少させたい})
  \end{cases}
\]
とおく．

\item \textbf{ピボット行の計算と比検定．} $\alpha^\top := e_r^\top A_B^{-1}A_N$(タブローの行 $r$ の非基底部分)を計算し，各非基底 $j\in N$ に対し $\sigma_j:=+1\,(\eta(j)=L)$，$\sigma_j:=-1\,(\eta(j)=U)$，$\hat\alpha_j:=\sigma_j\alpha_j$，$\hat c_j:=\sigma_j\bar c_j\,(\ge0)$ とおく．加入可能な非基底変数の集合
\[
  \mathrm{Eligible} \ :=\ \{\, j\in N : d\,\hat\alpha_j < 0 \,\}
\]
を計算する．$\mathrm{Eligible}=\emptyset$ ならば，現在の $M$ のもとで問題~\eqref{eq:lpM}は実行不可能であり，アルゴリズムを終了する．そうでなければ
\[
  q \ \in\ \operatorname*{arg\,min}_{j\in \mathrm{Eligible}} \frac{\hat c_j}{|\hat\alpha_j|}
\]
により加入変数 $q$ を選ぶ．

\item \textbf{ピボット．} $\beta(r)$ を非基底化し，$d=+1$ ならば $\hat l_{\beta(r)}$ に，$d=-1$ ならば $\hat u_{\beta(r)}$ に固定する．$q$ を基底化し，基底 $B$，逆行列 $A_B^{-1}$，基底解 $x_B$，非基底値 $x_N$，非基底の状態 $\eta$ を更新して(a)に戻る．
\end{enumerate}

\subsection{$M\to+\infty$ における挙動の解析}

Step~2(c)(d)(ピボット行・被約費用・加入可能集合・比検定・基底更新)は，現在の基底 $B$ と非基底の状態 $\eta$ のみから定まり，$\hat l,\hat u$ の数値(したがって $M$ の数値)には依存しない．$M$ が直接関与するのは，(i) 式~\eqref{eq:trunc}による $\hat l_j,\hat u_j$ 自体，(ii) それに応じて定まる $x_N,x_B$，(iii) Step~2(a)(b)の主実行可能性判定・脱退変数の選択・方向 $d$ の決定，の3箇所のみである．以下では，組合せ的状態 $(B,\eta)$ を固定したまま，これらが $M$ のアフィン関数(1次式)になることを用いて，$M\to+\infty$ でのアルゴリズムの挙動を解析する．

\begin{lemma}[$x_B$ の $M$ に関するアフィン性]
\label{lem:affine}
組合せ的状態 $(B,\eta)$ を固定する．非基底値 $x_N(M)$ の各成分は
\[
  x_j(M) =
  \begin{cases}
    l_j & (\eta(j)=L,\ l_j>-\infty)\\
    -M & (\eta(j)=L,\ l_j=-\infty)\\
    u_j & (\eta(j)=U,\ u_j<+\infty)\\
    M & (\eta(j)=U,\ u_j=+\infty)
  \end{cases}
\]
であるから，$x_N(M) = x_N^0 + M\delta$ と書ける．ここで $x_N^0\in\R^{|N|}$，$\delta\in\{-1,0,1\}^{|N|}$ は $(B,\eta)$ のみに依存し $M$ に依存しない定ベクトルで，$\delta_j=-1\,(\eta(j)=L,l_j=-\infty)$，$\delta_j=+1\,(\eta(j)=U,u_j=+\infty)$，それ以外は $\delta_j=0$ である．したがって
\[
  x_B(M) \ =\ A_B^{-1}\bigl(b-A_Nx_N(M)\bigr) \ =\ \underbrace{A_B^{-1}(b-A_Nx_N^0)}_{=:p}\ -\ M\underbrace{A_B^{-1}A_N\delta}_{=:\rho}
\]
は $M$ のアフィン関数であり，係数 $p,\rho\in\R^m$ は $(B,\eta)$ にのみ依存する(脱退行の添字として既に用いている $r$ との混同を避けるため，ここではベクトル係数を $\rho$ と表す)．
\end{lemma}

同様に，$\hat l_{\beta(i)}(M),\hat u_{\beta(i)}(M)$ も式~\eqref{eq:trunc}よりそれぞれ($l_{\beta(i)}$ または $-M$)，($u_{\beta(i)}$ または $M$)というアフィン関数であるから，Step~2(a)の逸脱量
\[
  v_i^-(M) := \hat l_{\beta(i)}(M) - x_{B,i}(M), \qquad
  v_i^+(M) := x_{B,i}(M) - \hat u_{\beta(i)}(M)
\]
も，補題~\ref{lem:affine}よりいずれも $M$ のアフィン関数である．これを $v_i^-(M)=a_i^-M+b_i^-$，$v_i^+(M)=a_i^+M+b_i^+$ と書く($a_i^\pm,b_i^\pm\in\R$ は $(B,\eta)$ のみに依存する定数)．

\begin{definition}[アフィン関数の辞書式順序]
2つのアフィン関数 $f(M)=aM+b$，$g(M)=cM+d$ に対し，
\[
  f \ \succ\ g \quad :\Longleftrightarrow\quad a>c,\ \ \text{または}\ \ (a=c\ \text{かつ}\ b>d)
\]
と定める(すなわち係数ベクトル $(a,b),(c,d)$ の辞書式順序で比較する)．
\end{definition}

$f\succ g$ ならば，ある $M^\dagger(f,g)>0$ が存在し，すべての $M\ge M^\dagger(f,g)$ で $f(M)>g(M)$ が成り立つ($a\neq c$ なら $M^\dagger=|b-d|/|a-c|$ でよく，$a=c,\,b>d$ なら任意の $M>0$ で成り立つ)．$g\succ f$ の場合も同様であり，$(a,b)=(c,d)$ のときは $f\equiv g$ である．この $\succ$ は有限個のアフィン関数の集合上の全順序を与え，それらすべての対に対する $M^\dagger(f,g)$ の最大値を $M^\dagger$ とおけば，$M\ge M^\dagger$ である限り，数値比較の結果はすべて $\succ$ による比較と一致する．

これを用いて，Step~2(a)(b)を，具体的な $M$ の値を固定せずに次のように書き換えることができる．

\begin{enumerate}
\item[(a$'$)] \textbf{主実行可能性の判定(終了条件，$M\to\infty$ 版)．} $V_\infty := \{\, i\in\{1,\dots,m\} : v_i^- \succ 0 \ \text{または}\ v_i^+ \succ 0 \,\}$ とおく(ここで $0$ は恒等的に $0$ であるアフィン関数，すなわち係数ペア $(0,0)$ を表す)．$V_\infty=\emptyset$ ならば，$M\ge M^\dagger$ のすべてで主実行可能性が成り立ち，アルゴリズムは終了する．
\item[(b$'$)] \textbf{脱退変数の選択($M\to\infty$ 版Dantzig則)．} 各 $i\in V_\infty$ に対し，$v_i^-,v_i^+$ のうち $\succ$ で大きい方を $w_i$ とする．
\[
  r \ \in\ \operatorname*{arg\,max}_{i\in V_\infty}{}^{\succ}\, w_i
\]
により(辞書式順序 $\succ$ に関する最大元として)脱退行 $r$ を選ぶ．$w_r=v_r^-$ であれば $d=+1$，$w_r=v_r^+$ であれば $d=-1$ とする(両者が等しい場合の同点処理は，添字の小さい方を優先するなど，Bland 則に準じた規則で決定的に定める)．
\end{enumerate}

\begin{proposition}[数値計算との一致]
組合せ的状態 $(B,\eta)$ を固定したとき，(a$'$)(b$'$)による判定・選択は，ある(その状態にのみ依存する)$M^\dagger>0$ が存在して，すべての $M\ge M^\dagger$ における(a)(b)の数値計算による判定・選択と一致する．
\end{proposition}

\begin{proof}
$V_\infty$ の各元の判定，および $\arg\max^\succ$ による $r$ の選択は，有限個のアフィン関数の対 $(v_i^-,v_i^+,\,0)$ の間の $\succ$ による比較のみから決まる．上述の通り，これらの比較はある共通の $M^\dagger$ 以降，数値比較 $v_i^\pm(M)>0$，$\max_i(\cdot)$ と一致する．
\end{proof}

Step~2(c)(d)は $M$ に依存しないから，(a$'$)(b$'$)によって次の状態 $(B',\eta')$ が(数値としての $M$ を固定することなく)一意に定まり，これを繰り返すことで，組合せ的な状態遷移の列 $(B^{(0)},\eta^{(0)}) \to (B^{(1)},\eta^{(1)}) \to \cdots$ が得られる．この列は，可能な基底の総数が有限であることと(退化に対する)巡回防止規則(今後の課題とする)により有限回で停止し，停止する反復回数を $K$ とすれば，各反復で用いた $M^\dagger$ の最大値を改めて $M^\dagger$ とおくことで，$M\ge M^\dagger$ なるすべての(数値としての)$M$ に対して，Step~2を実際に数値計算した場合の反復列が，この組合せ的な遷移列と完全に一致することがわかる．すなわち，$M\to+\infty$ での挙動は，具体的な $M$ の値を固定せずに，係数ペア $(a,b)$ の辞書式比較のみによって先取りして計算できる．

\paragraph{最適値の $M$ に関するアフィン性と非有界性の判定．}
状態遷移列が最終状態 $(B^\ast,\eta^\ast)$((a$'$)で $V_\infty=\emptyset$ となる状態)に到達したとする．補題~\ref{lem:affine}の $p,\rho$($x_B(M)=p-M\rho$)と $x_N(M)=x_N^0+M\delta$ を用いると，問題~\eqref{eq:lpM}の最適値は
\[
  z(M) \ =\ c_B^\top x_B(M) + c_N^\top x_N(M)
     \ =\ \underbrace{\bigl(c_B^\top p + c_N^\top x_N^0\bigr)}_{=:z_0} \ +\ M\ \underbrace{\bigl(c_N^\top\delta - c_B^\top \rho\bigr)}_{=:z_1}
\]
と，やはり $M$ のアフィン関数として表される．ここで $z_0,z_1\in\R$ は最終状態 $(B^\ast,\eta^\ast)$ にのみ依存する定数であり，$M\ge M^\dagger$ の範囲では，この $z(M)=z_0+z_1M$ が問題~\eqref{eq:lpM}の真の最適値そのものである(状態が $(B^\ast,\eta^\ast)$ のまま変化しないため)．

式~\eqref{eq:trunc}より $M_1\le M_2$ ならば $[\hat l(M_1),\hat u(M_1)]\subseteq[\hat l(M_2),\hat u(M_2)]$(有限であった側は不変，非有界であった側のみ広がる)であるから，\eqref{eq:lpM}の実行可能領域は $M$ に関して単調非減少であり，最小値である $z(M)$ は $M$ に関して単調非増加でなければならない．アフィン関数 $z(M)=z_0+z_1M$ がこの性質を持つのは $z_1\le0$ の場合に限られるから，
\[
  z_1 \ \le\ 0
\]
が常に成り立つ．

\begin{lemma}[実行可能領域の合併]
\label{lem:union}
$F:=\{x:Ax=b,\ l\le x\le u\}$(元問題~\eqref{eq:lp}の実行可能領域)とする．このとき
\[
  F \ =\ \bigcup_{M>0} F_M, \qquad F_M := \{x:Ax=b,\ \hat l(M)\le x\le \hat u(M)\}\ (\text{すなわち\eqref{eq:lpM}の実行可能領域})
\]
が成り立つ．
\end{lemma}

\begin{proof}
($\supseteq$)：任意の $M>0$ に対し $[\hat l(M),\hat u(M)]\subseteq[l,u]$(式~\eqref{eq:trunc}，有限であった側は不変，非有界であった側のみ真に拡張される)であるから $F_M\subseteq F$．ゆえに $\bigcup_M F_M\subseteq F$．

($\subseteq$)：$x\in F$ を任意にとる．$x$ は(拡大実数値ではなく)有限な実ベクトルであるから $\|x\|_\infty<\infty$．$M\ge\|x\|_\infty$ ととれば，$x_j$ がもともと有限な上下限を持つ座標では $\hat l_j(M)=l_j\le x_j\le u_j=\hat u_j(M)$ がそのまま成り立ち，もともと非有界であった座標では $|x_j|\le M$ より $-M\le x_j\le M$ が成り立つ．ゆえに $x\in F_M$．
\end{proof}

この符号によって，$M\to\infty$ での挙動，したがって元問題~\eqref{eq:lp}の性質が次のように読み取れる．

\begin{itemize}
\item $z_1<0$ のとき：$z(M)\to-\infty\ (M\to\infty)$ である．$[\hat l(M),\hat u(M)]\subseteq[l,u]$(第4.2節)より \eqref{eq:lpM}の実行可能領域は元問題の実行可能領域 $F:=\{x:Ax=b,\ l\le x\le u\}$ の部分集合であるから，$\inf_{x\in F}c^\top x \ \le\ z(M) \ \to\ -\infty$ となり，元問題は(実行可能かつ)\emph{非有界}であると判定できる．
\item $z_1=0$ のとき：$z(M)=z_0$ は $M\ge M^\dagger$ で一定であり，人工上下限を有限な $M$ からさらに緩めても最適値が変化しないことを意味する．補題~\ref{lem:affine}より，最終状態 $(B^\ast,\eta^\ast)$ における非基底値は $x_N(M)=x_N^0+M\delta$ という形をしていた．各成分ごとに書けば，
\[
  x_j(M) =
  \begin{cases}
    l_j\ \text{または}\ u_j\ (\text{その真の値，一般には非零}) & (\delta_j=0,\ \text{すなわちその側がもともと有限}) \\
    \mp M & (\delta_j=\mp1,\ \text{すなわちその側がもともと非有界})
  \end{cases}
\]
である．両側とも有限であった(真に有界な)変数は $\delta_j=0$ であり，その非基底値は $M$ によらず常に真の(一般には非零の)上下限のままである点に注意されたい($0$ になるのは，第4.1節の前処理によって有限側がちょうど $0$ に正規化された，片側のみ非有界であった変数の場合に限られる)．形式的に $M=0$ を代入すると，$\delta_j=0$ の成分は真の有限値のまま変わらず，$\delta_j=\mp1$ の成分(元々非有界であった側に配置されていた非基底変数)のみが $\mp M\to0$ となる．いずれの場合も，これはちょうど $x_N^0$ の定義(補題~\ref{lem:affine})に一致するから，$x_N(0)=x_N^0$ である．したがって
\[
  x^\ast \ :=\ \bigl(x_B(0),\,x_N(0)\bigr) \ =\ (p,\,x_N^0)
\]
は前処理後の元問題~\eqref{eq:lp}の実行可能解であり，$c^\top x^\ast = z(0) = z_0$(アフィン関数 $z(M)=z_0+z_1M$ に $z_1=0$ を用いれば直ちに $z_0$ に一致する)を満たす，具体的な最適解を与える．すなわち $z_1=0$ の場合，$M\to+\infty$ の極限を考える必要はなく，形式的に $M=0$ とおいた(切断を全く行わない)場合の最終状態の基底解が，そのまま元問題の最適解になる．
\end{itemize}

すなわち，最適値の $M$ に関する傾き $z_1$ の符号($<0$ か $=0$ か)だけから，具体的な $M$ の値によらず，元問題が非有界であるか有限な最適値を持つかを判定でき，$z_1=0$ の場合にはその最適解の具体的な値まで($M=0$ の代入により)直ちに得られる．この判定を実行可能性の判定(前述の $V_\infty=\emptyset$ となるか否か)と組み合わせれば，元問題の実行不可能性・非有界性・有限最適解の存在という3通りの場合を過不足なく区別できることが，次の命題として実際に証明できる．

\begin{proposition}[3通りの場合の判別]
\label{prop:trichotomy}
組合せ的状態遷移列が，ある $M^\dagger>0$ 以降で最終状態 $(B^\ast,\eta^\ast)$ に到達し，そこで(a$'$)(b$'$)(c)(d)が停止するとする．このとき次が成り立つ．
\begin{enumerate}[label=(\roman*)]
\item 停止結果が \textsc{Infeasible} である $\iff$ $F=\emptyset$．
\item 停止結果が \textsc{Optimal} かつ $z_1<0$ である $\iff$ $F\neq\emptyset$ かつ $\inf_{x\in F}c^\top x=-\infty$(元問題は実行可能かつ非有界)．
\item 停止結果が \textsc{Optimal} かつ $z_1=0$ である $\iff$ $F\neq\emptyset$ かつ $\inf_{x\in F}c^\top x$ が有限であり，このとき $x^\ast=(p,x_N^0)$ がその最適解，$z_0$ がその最適値を与える．
\end{enumerate}
\end{proposition}

\begin{proof}
第4.5節冒頭で示した通り，(a$'$)(b$'$)(c)(d)による組合せ的な状態遷移は，$M\ge M^\dagger$ のすべての(数値としての)$M$ に対し，Step~2を実際に数値計算した場合の反復と完全に一致する．したがって停止結果は $M\ge M^\dagger$ のすべてで共通であり，$F_M$ もまた $M\ge M^\dagger$ のすべてで共通(空集合であるか，あるいは常に非空で最小値 $z(M)=z_0+z_1M$ を持つか)である．また，式~\eqref{eq:trunc}より $M_1\le M_2\Rightarrow[\hat l(M_1),\hat u(M_1)]\subseteq[\hat l(M_2),\hat u(M_2)]$，すなわち $F_{M_1}\subseteq F_{M_2}$(単調性)である．

(i)：($\Leftarrow$) $F=\emptyset$ ならば，補題~\ref{lem:union}より任意の $M$ で $F_M\subseteq F=\emptyset$．とくに $M\ge M^\dagger$ で \eqref{eq:lpM} は実行不可能であり，停止結果は \textsc{Infeasible} である．
($\Rightarrow$) 停止結果が \textsc{Infeasible} であれば $F_{M^\dagger}=\emptyset$．単調性より $M<M^\dagger$ でも $F_M\subseteq F_{M^\dagger}=\emptyset$．ゆえに補題~\ref{lem:union}より $F=\bigcup_M F_M=\emptyset$．

(ii)(iii)：以下，(i)より $F\neq\emptyset$ であるとしてよく，したがって停止結果は \textsc{Optimal} である．補題~\ref{lem:union}と単調性より，任意の $x\in F$ に対し $M_x:=\max(M^\dagger,\|x\|_\infty)$ とおけば $x\in F_{M_x}$(補題~\ref{lem:union}の証明と同じ議論)であり，ゆえに
\[
  c^\top x \ \ge\ \min_{F_{M_x}} c^\top x \ =\ z(M_x) \ =\ z_0+z_1M_x .
  \tag{$\ast$}
\]

($z_1<0 \Rightarrow$ 非有界)：$z_1<0$ のとき，$(\ast)$ で $x$ を固定したまま $M_x\to\infty$ ととる必要はなく，むしろ $z(M)\to-\infty\,(M\to\infty)$ かつ $z(M)=\min_{F_M}c^\top x\ge\inf_{x\in F}c^\top x$(単調性より $F_M\subseteq F$)であることから，$\inf_{x\in F}c^\top x\le\lim_{M\to\infty}z(M)=-\infty$．

(非有界 $\Rightarrow z_1<0$)：対偶を示す．$z_1=0$ とする．$F\neq\emptyset$ かつ非有界であれば，任意の $N>0$ に対しある $x_N\in F$ が存在して $c^\top x_N<-N$．$(\ast)$ を $x=x_N$ に適用すると $c^\top x_N\ge z_0+z_1M_{x_N}=z_0$($z_1=0$ のため $M_{x_N}$ の値によらない)．ゆえに $z_0\le c^\top x_N<-N$．$N$ は任意だから $z_0=-\infty$ となるが，$z_0$ は補題~\ref{lem:affine}より有限な実数として定義されているから矛盾．ゆえに $F$ が非有界ならば $z_1\neq0$．先に示した $z_1\le0$ と合わせて $z_1<0$．

以上より，$F\neq\emptyset$ のとき「$z_1<0\iff$ 非有界」，したがって($z_1\le0$ の二分法により)「$z_1=0\iff$ 有限な最適値を持つ」が従う．

最後に，$z_1=0$ の場合の最適性を示す．任意の $x\in F$ に対し，$(\ast)$ で $z_1=0$ を用いれば $c^\top x\ge z_0$．一方，先に示した通り $x^\ast=(p,x_N^0)\in F$ かつ $c^\top x^\ast=z_0$．ゆえに $z_0=\min_{x\in F}c^\top x=\inf_{x\in F}c^\top x$ であり，$x^\ast$ はその最適解である．
\end{proof}

\begin{remark}
最終状態(a$'$)で $V_\infty=\emptyset$ となったときの最適解において，非基底変数 $j$ の値は常に厳密に $x_j=\hat l_j$ または $\hat u_j$ であるから，元々 $l_j=-\infty$(または $u_j=+\infty$)であった $j$ が人工上下限に拘束されているかどうかは，最終的な非基底の状態 $\eta$ が $\eta(j)=L$(または $U$)であるかどうかという\emph{組合せ的な}条件に帰着し，$M$ の数値を代入して確認する必要はない．
\end{remark}

\paragraph{最適基底の特定．}
上記の $z_1,z_0,p,x_N^0$ はいずれも，最終状態 $(B^\ast,\eta^\ast)$——(a$'$)(b$'$)とStep~2(c)(d)を，数値としての $M$ を固定せずに反復した，組合せ的な状態遷移列の終着点——にのみ依存する．したがって，最適基底 $(B^\ast,\eta^\ast)$ を特定するために，何らかの具体的な数値を $M$ に代入して単体法を実行し直す必要はない．(a$'$)(b$'$)による判定・選択(および最急辺規則・Devex規則を用いる場合は後述の係数3つ組の辞書式比較)と，$M$ に依存しないStep~2(c)(d)とを繰り返し，反復が停止した時点の状態がそのまま最適基底 $(B^\ast,\eta^\ast)$ である．とくに，停止後に $z_1$ を計算して $z_1=0$ であることが確認できれば，その $(B^\ast,\eta^\ast)$ における $p,x_N^0$ をそのまま用いて，上で述べた $x^\ast=(p,x_N^0)$ が元問題の具体的な最適解として得られる．

\paragraph{非基底の自由変数の扱い．}
最終状態 $(B^\ast,\eta^\ast)$ で非基底のまま残る変数 $j$ の中に，片側だけでなく\emph{両側とも}非有界な自由変数($l_j=-\infty$ かつ $u_j=+\infty$)が含まれる場合を考える．その値は $\eta(j)=L$ なら $x_j(M)=-M$，$\eta(j)=U$ なら $x_j(M)=M$ であり，$z_1=0$ の場合には $M=0$ の代入で $x_j^\ast=0$ としていた．この値の選び方が妥当であることは，次の観察から正当化される．

標準的な単体法の議論により，$x_B(M)=A_B^{-1}(b-A_Nx_N(M))$ を用いて $z(M)=c_B^\top x_B(M)+c_N^\top x_N(M)$ を書き直すと，
\[
  z(M) \ =\ c_B^\top A_B^{-1}b \ +\ \sum_{j\in N} \bar c_j\, x_j(M)
\]
となる．非基底値 $x_j(M)$ のうち $M$ に依存する項(人工上下限に置かれた $j$)だけを取り出せば，$z(M)=z_0+z_1M$ の傾きは
\[
  z_1 \ =\ \sum_{\substack{j\in N,\ \eta(j)=L\\ l_j=-\infty}} (-\bar c_j) \ +\ \sum_{\substack{j\in N,\ \eta(j)=U\\ u_j=+\infty}} \bar c_j
\]
と表される(この式は，補題~\ref{lem:affine}の $z_1=c_N^\top\delta-c_B^\top\rho$ を被約費用で書き直したものに等しい)．双対実行可能性より，$\eta(j)=L$ の各項は $\bar c_j\ge0$ ゆえ $-\bar c_j\le0$，$\eta(j)=U$ の各項は $\bar c_j\le0$ であるから，$z_1$ は\emph{非正の項の和}であり，これは先に示した $z_1\le0$ の別証明でもある．とくに，
\[
  z_1=0 \quad\Longleftrightarrow\quad \text{人工上下限($\mp M$)に置かれたすべての非基底変数 $j$ について}\ \bar c_j=0
\]
が成り立つ．

これはまさに，非有界な方向に対する標準的な最適性条件——非有界な方向に置かれた非基底変数の被約費用はちょうど $0$ でなければならない(さもなくばその方向に動かすことで目的関数をさらに減少させられ，実際 $\bar c_j\neq0$ であれば $z_1\neq0$ となり(ii)の非有界の場合に該当する)——に他ならない．とくに自由変数 $j$ がこの条件を満たす場合，$\bar c_j=0$ は，他の基底変数を $Ax=b$ を保つように追随させながら $x_j$ をどの値に動かしても目的関数の値が変化しないことを意味する．すなわち，$x_j^\ast=0$ という値そのものに特別な意味はなく，$x_N^0$ の定義上その座標がたまたま $0$ になるという便宜的な代表点に過ぎない．元問題の最適解は，一般にはこの自由変数の値に関して一意ではなく(($\bar c_j=0$ の非基底変数を動かして得られる)同一目的関数値を持つ面が存在し得る)，$x^\ast=(p,x_N^0)$ はその面上の一つの代表点を与えるものと理解すべきである．一方，最適\emph{基底} $(B^\ast,\eta^\ast)$ 自体の特定は，自由変数の存在によって何ら妨げられない．

\begin{remark}[自由変数を非基底のままとしてよいか]
最終状態で自由変数 $j$ が非基底のまま残ること自体は，古典的な(上下限のない)単体法の教科書的な扱い——自由変数は常に基底に含める，という規約——とは一見異なるように見えるが，矛盾するものではない．むしろ，$z_1=0$ のとき $\bar c_j=0$ であることは，この自由変数を非基底のままにしておいても\emph{最適性が損なわれない}ことの証拠そのものである．そして，$\bar c_j=0$ は同時に，この $j$ を(どこかの基底変数と交換して)基底に組み入れる\emph{退化的な}ピボット(目的関数値を変化させないピボット)を行うことも常に可能であることを意味する．すなわち，古典的な規約に合わせて「自由変数は必ず基底に含める」形の最適基底を望むならば，$\bar c_j=0$ の非基底変数 $j$ を加入変数とする追加のピボットを行えばよく，これによって同じ最適値 $z_0$ を保ったまま，自由変数が基底に含まれる形の(同値な)最適基底へ移ることができる．逆に言えば，本稿の枠組みでは，そのような追加のピボットを行う必要は\emph{ない}――非基底のままの自由変数を許した形の最適基底で，最適値の判定という目的には既に十分である．
\end{remark}

\paragraph{最急辺規則・Devex規則への拡張．}
Dantzig則((b$'$))はアフィン関数どうしの比較(前述の $\succ$，次数1)に帰着したが，最急辺規則・Devex規則は逸脱量を重みで正規化した $\Delta_i^2/\gamma_i$(あるいは $\Delta_i^2/w_i$)を比較するため，一般には $M$ の2次式どうしの比較になる．まず，重み自体が $M$ に依存しないことを確認する．

\begin{lemma}[重みの $M$ 独立性]
最急辺規則の重み $\gamma_i=\sum_{j\in N}\alpha_{ij}^2$ は，タブロー行の成分 $\alpha_{ij}=(A_B^{-1}A_N)_{ij}$ のみから定まり，Step~2(c)で述べた通り $\alpha_{ij}$ は $A$ と現在の基底 $B$ のみに依存して $\hat l,\hat u$ の数値(したがって $M$)には依存しない．Devex規則の重み $w_i$ も，初期値 $1$ とピボット行の成分 $\alpha_{rj}$ のみによる更新から定まるため，同様に $M$ に依存しない．
\end{lemma}

したがって，(b$'$)で定めた脱退候補行 $i\in V_\infty$ の逸脱量 $w_i(M)=a_iM+b_i$(組合せ的状態を固定したときの，第4.5節冒頭の記法)を用いると，
\[
  \frac{w_i(M)^2}{\gamma_i} \ =\ \frac{a_i^2}{\gamma_i}\,M^2 \ +\ \frac{2a_ib_i}{\gamma_i}\,M \ +\ \frac{b_i^2}{\gamma_i}
\]
は $M$ の高々2次式であり($\gamma_i$ は $M$ に依存しないから)，Devex規則についても $\gamma_i$ を $w_i$ に置き換えれば同様である．これらを比較するために，アフィン関数の辞書式順序を一般の多項式へ拡張する．

\begin{definition}[多項式の辞書式順序]
次数 $d$ 以下の多項式 $f(M)=\sum_{k=0}^{d}a_kM^k$，$g(M)=\sum_{k=0}^{d}c_kM^k$ に対し，$k^*:=\max\{\,k : a_k\neq c_k\,\}$ が存在するとき
\[
  f \ \succ\ g \quad :\Longleftrightarrow\quad a_{k^*} > c_{k^*}
\]
と定める(存在しなければ $f\equiv g$)．先のアフィン関数の辞書式順序は $d=1$ の場合にあたる．
\end{definition}

多項式の最高次の非零係数がその符号を(十分大きな $M$ で)決定するという初等的な性質により，$f\succ g$ ならばある $M^\dagger(f,g)>0$ が存在して $M\ge M^\dagger(f,g)$ で $f(M)>g(M)$ となる．これにより，最急辺規則・Devex規則による脱退変数選択の $M\to\infty$ 版は，3つ組(次数2,1,0の係数)
\[
  \Bigl(\tfrac{a_i^2}{\gamma_i},\ \tfrac{2a_ib_i}{\gamma_i},\ \tfrac{b_i^2}{\gamma_i}\Bigr)
\]
の辞書式比較として
\[
  r \ \in\ \operatorname*{arg\,max}_{i\in V_\infty}{}^{\succ}\ \frac{w_i(M)^2}{\gamma_i}
\]
により，Dantzig則の(b$'$)と全く同様に，数値としての $M$ を固定せずに先取りして計算できる．

\begin{remark}
Dantzig則は $\gamma_i\equiv1$ とした特別な場合ではなく，比較する量の次数そのものが異なる($w_i$ 自体か $w_i^2/\gamma_i$ か)点に注意されたい．しかし，係数ベクトルの辞書式比較という解析の枠組みは3つの規則に共通しており，比較する量の次数 $d$ に応じて，比較すべき係数ベクトルの長さ($d+1$ 成分)が変わるだけである．
\end{remark}

\end{document}
